\section*{AI-use statement}

This paper was drafted by a crew of AI research-assistant agents
(Free-Research-Crew, protocol v0.1) under the direction of the human author, Bryce Wang.
He approved the research question (2026-11-19, 01:17 UTC), the design (2026-11-19, 07:15 UTC)
and the results (2026-11-19, 12:27 UTC) at three human gates, each in chat after reading the
Reviewer's verdict and kill argument. The manuscript itself is subject to a fourth, the
submission gate. He takes responsibility for the paper's content.

\begin{itemize}
  \item \textbf{Roles and models.} Every crew role ran on Anthropic Claude through Claude
        Code. The PI role was the main Claude Code session. The Scout, Designer, Data RA,
        Experimenter, Figures RA, Reviewer and Writer ran as agents launched from that
        session; the review records name the model as \texttt{claude-opus-5-5}. The PI's
        brief to each agent is archived in \texttt{briefs/}, from the Scout's first brief
        (2026-11-18, 16:28 UTC) to this Writer's brief.
  \item \textbf{Review.} The Reviewer ran 3 rounds on the research question, 2 on the
        design, 1 on the code and 2 on the results; the paper review comes before the
        submission gate, and this item is to be updated after it. No separate literature
        round was held. The code review passed with no critical finding before any analysis
        run. The Reviewer's independence level was \emph{fresh-context}: a new context in
        every round, with the same model family as the other roles. The human author chose
        this over a cross-family review. The last round of each stage passed with minor
        findings still open: question round~3 (deferred to the design and answered there),
        design round~2 (one, answered and disclosed before the design gate), code round~1
        (verdict \emph{revise}, no critical finding; one fixed and two answered by narrowing
        the plan before any analysis run, Section~\ref{sec:deviations}) and results round~2
        (three, answered in the run notes and claims before the results gate).
  \item \textbf{Findings accepted over the Reviewer's objection.} None. Where the crew
        narrowed a claim instead of changing the analysis (the H1 wording at design round~1
        and code round~1), the Reviewer recorded the finding as conceded.
  \item \textbf{Evidence.} Every result number in the text and tables is a macro generated
        from the metrics of a logged run, or is printed by a logged table script from run
        metrics. Numbers typed by hand are design facts fixed in the plan before any
        analysis run (dataset sizes and positive shares, split sizes, seeds, margins, the
        clipping constant, the bin count, the noise-floor approximation), facts from the
        cited sources as the Scout's notes record them, and process facts such as dates and
        round counts. Every citation was checked against its source by the Scout: of the 23
        sources, 15 were read in full text, 6 from the abstract only and 2 by resolving their
        DOI, and each is cited for no more than was read. The full workspace is public at
        \url{https://github.com/brycewang-stanford/Free-Research-Crew/tree/main/gallery/cs-calibration}.
        It contains the plan with its disclosures, the run ledger and notes, the code, the
        data manifest, every review, the archived briefs and the source notes.
  \item \textbf{Replication.} Pending. The clean rerun of every logged run from
        \texttt{code/requirements.txt} (task \texttt{r-rerun}) has not yet been done, and no
        \texttt{replication.json} exists; this item is to be updated when it does.
\end{itemize}
